\chapter*{Abstract}
\addcontentsline{toc}{chapter}{Abstract}

The standard presentation of relativistic quantum mechanics relies on matrix
representations and complex vector spaces that obscure the geometric
foundations of the theory.
This monograph investigates the origins of the Dirac theory, tracing its
emergence from fundamental relativistic principles into the coordinate-free
framework of real Clifford algebras.
Because the invariance of the speed of light defines a pseudo-Euclidean
quadratic space, the spacetime algebra $\mathcal{C}\ell_{1,3}$ arises naturally
by internalising the Minkowski metric into an associative product.
Within this algebraic setting, the linearisation of the second-order
Klein--Gordon--Fock equation is shown to be a direct consequence of the geometry
of quadratic spaces.
The Dirac spinor is identified as an even multivector, replacing abstract
complex matrix spaces with real spacetime geometry and interpreting the
imaginary unit via spatial bivectors and the spacetime pseudoscalar.
Intrinsic spin is shown to originate from the rotational symmetries and covering
properties of the spin group rather than as an ad hoc postulate.
The electromagnetic gauge coupling, the non-relativistic Pauli--Schrödinger limit
with the Darwin correction, the free normal modes, and the discrete $C$, $P$,
and $T$ symmetries are examined.
This analysis clarifies the geometric foundation underlying the Dirac equation,
providing a unified mathematical description that connects relativistic
mechanics with the algebraic structures required for continuous matter fields.
